% Printed 498; principal PDF 505; Google PDF 509.
\par\clearpage
\noindent\makebox[0pt][l]{\fontsize{13}{16}\selectfont 498}\hfill{\small\textbf{BRIEFWECHSEL.}}\hfill\null\par\medskip
\noindent dass in meiner Sternwarte keine recht schickliche Aufstellung zu erhalten ist zu dergleichen Probemessungen, auch wenig geeignete Objecte sichtbar sind, endlich die Winterjahreszeit für meinen sehr empfindlichen Körper ungünstig ist. Ich will aber wenigstens versuchen, was in meinen Kräften steht.

\par\noindent\hspace*{\parindent}Sollten Sie nun vielleicht selbst ähnliche Erfahrungen, und in deren Veranlassung allerlei darauf bezügliche Versuche gemacht haben, so werden Sie mich durch Mittheilung sehr verpflichten. Es ist in der That eine Lebensfrage bei Theodolithen von dieser Construction. Bei \textsc{Borda}schen Kreisen kann der Fehler vermöge des Constructionsprincips gar nicht eintreten; es ist gar kein Grund da, dass Bewegung der Alhidade die Lage des untern Fernrohrs \textls[120]{gegen den Limbus}kreis verrücke, und ebenso wenig verrückt die Bewegung des ganzen Kreises die Lage der Alhidade gegen den Limbuskreis. Man kann Theodolithen nach demselben Princip bauen, aber dann sind es ganz andere Instrumente wie die \textsc{Reichenbach}schen; wenn ich nicht irre, hat \textsc{Schumacher} einen solchen von \textsc{Gambey} (wo dann das untere Fernrohr eine ganz andere Rolle spielt); ich kenne aber das Detail der Einrichtung nicht.\enspace .\ .\ .\ .\ .\ .

\par\noindent\hspace*{\parindent}Es ist auffallend, wie Sachen zu Papiere gebracht einen andern Eindruck machen, als wenn man sie nur im Kopfe überdenkt. Indem ich obige Zeilen noch einmal flüchtig übersehe, kommt es mir vor, dass einem bloss unbefangenen Leser, der selbst gar keine eigene Erfahrungen gemacht hatte, ein Umstand importanter erscheinen muss, als ich ihn selbst bisher gehalten habe, der Umstand, dass alle Theodolithen \textls[120]{zu Anfang} diese Erscheinung gar nicht oder nicht merklich gezeigt haben. Irgend eine \textls[120]{Ausnutzung} muss also nothwendig im Spiel sein, und ich werde also doch wohl jedenfalls neue Schrauben, Muttern und Kugeln machen lassen, obgleich, wie gesagt, irgend eine Vacillation an diesen Theilen im Einzelnen nicht erkannt werden kann. Ich werde aber versuchen, ob vielleicht ein \textls[120]{viel stärkeres} Fernrohr doch im Einzelnen schon etwas von einer solchen Vacillation erkennen lassen wird.

\par\bigskip
\begin{center}{\large\textsc{Gauss} an \textsc{Bessel}.\quad Göttingen, 15. August 1844.}\end{center}
\par\noindent\hspace*{\parindent}.\ .\ .\ .\ .\ .\enspace Der Gegenstand, worüber ich mit Ihnen zuletzt correspondirte, nemlich die Tendenz der Repetitions-\textls[120]{Theodolithen}, die Winkel zu klein zu

% Printed 499; principal PDF 506; Google PDF 510.
\par\clearpage
\noindent\null\hfill{\small\textbf{MESSUNGSFEHLER.}}\hfill\makebox[0pt][r]{\fontsize{13}{16}\selectfont 499}\par\medskip
\noindent geben, hat mich in vorigem Winter noch sehr geplagt. Die Ursache kann keine andere sein, als dass der Kreis nicht absolut fest bleibt während derjenigen Manipulationen der Alhidade, bei welchen vorausgesetzt wird, dass jener fest bleibe. Die Abänderungen aber, die ich habe anbringen lassen, haben sich sehr wirksam bewiesen, nicht nur während der äusserst zahlreichen Probemessungen, die ich selbst im Januar bis März machte, sondern auch bei dem wirklichen Gebrauch, den mein Sohn noch fortwährend von diesem Instrument macht. Die Resultate des vorigen Jahres waren zum Theil so schlecht, dass sie ganz verworfen werden mussten, während sie dies Jahr sämmtlich so gut sind, wie man von einem Instrument von dieser Dimension nur erwarten kann. Die Hauptveränderung besteht übrigens darin, dass Alhidadenzapfen und dessen Büchse (die ihrerseits den Limbuskreiszapfen bildet) jedes für sich durch eine Tragfeder dergestalt unterstützt wird, dass man jede dieser Tragfedern unabhängig von der andern nach Gefallen und augenblicklich anspannen oder abspannen (ganz oder theilweise) kann. Bei dem Gebrauch findet dann ein \textls[120]{planmässiger} beständiger Wechsel zwischen den Zuständen dieser Federn statt, so dass, alles gezählt, für jede Repetition 19 Operationen erfordert werden (das Ablesen mitgezählt); indessen macht man sich die Reihenfolge dieser Operationen so mechanisch, dass man fast ebenso schnell operiren kann, wie bei der gewöhnlichen Einrichtung.
\par\bigskip
\begin{center}\rule{2.6cm}{.4pt}\end{center}
\vspace{1.2cm}
{\small
\begin{center}BEMERKUNGEN.\end{center}
\par\noindent\hspace*{\parindent}Die Notiz über die bei der hannoverschen Gradmessung und Landesvermessung benutzten sowie für die letztere noch erforderlichen Instrumente ist einem im Gauss-Archiv befindlichen Bericht von \textsc{Gauss} vom 21. März 1829 an den Geheimen Cabinetsrath \textsc{Hoppenstedt} entnommen.
\par\noindent\hspace*{\parindent}Zu den vorstehenden Briefen an \textsc{Olbers}, \textsc{Schumacher} und \textsc{Bessel}, die nach den Originalen abgedruckt sind, ist die Stationsausgleichung für Brillit, S. 265/267, sowie auch eine Bemerkung von Professor \textsc{Goldschmidt}, S. 289, zu vergleichen.
\par\noindent\hfill\textsc{Krüger.}\hspace*{1em}\par}
\vspace{1cm}
\begin{center}\rule{3.5cm}{.4pt}\end{center}
\par\medskip\noindent\makebox[.86\linewidth][r]{63\textsuperscript{*}}\par

% Printed 500; principal PDF 507; Google PDF 511.
\par\clearpage
\noindent\makebox[0pt][l]{\fontsize{13}{16}\selectfont 500}\hfill{\small\textbf{BRIEFWECHSEL.}}\hfill\null\par
\vspace*{3.2cm}
\begin{center}
{\LARGE BRIEFWECHSEL UND NACHLASS.}\par
\vspace{1.5cm}
{\large [Berechnung des mittlern Ablesungsfehlers einschliesslich des Theilungsfehlers\par am Theodolithen.]}\par
\bigskip\bigskip
{\large [1.]}\par\medskip
{\large\textsc{Gauss} an \textsc{Gerling}.\quad Göttingen, 17. April 1844.}
\end{center}
\par\noindent\hspace*{\parindent}.\ .\ .\ .\ .\ .\enspace Über eine andere Prüfungsrechnung an Theodolithenbeobachtungen (welche zum Zweck hat, den mittlern Fehler der \textls[120]{Ablesungen} incl. der Theilungsfehler zu bestimmen) muss ich mir vorbehalten, Ihnen ein andermal zu schreiben. Ich habe diese Rechnungen schon vor 19 Jahren ausgeführt, und sie haben mir viel Vergnügen gemacht. Ich habe nemlich für \textls[120]{jede} vorgekommene vollständige Ablesung $a-b+c-d=l$ berechnet, wo (mit Weglassung der Grade) $a$, $b$, $c$, $d$ die 4 Ablesungen bedeuten. Dies $l$ sollte (\textls[120]{trotz} der Excentricität, ja trotz einer etwa veränderlichen Excentricität) \textls[120]{constant} sein; es ist nur variabel in Folge der Ablesungs- und der Theilungsfehler. Setzt man den mittlern Werth aus sehr vielen, $n$, Ablesungen $=\lambda$ (der bei einem gegebenen Theodolithen aus vielen Beobachtungen sich sehr scharf finden lässt und unveränderlich sein sollte) und $\frac{\Sigma(l-\lambda)^2}{n-1}=kk$, so ist $\frac12k$ der mittlere Fehler, den man bei der Ablesung Eines Index, und $\frac14k$ der mittlere Fehler, den man bei dem Mittel aus allen 4 Indicibus riskirt. Ich finde nun den letztern oder $\frac14k$ aus vollständigen Ablesungen:

% Printed 501; principal PDF 508; Google PDF 512.
\par\clearpage
\noindent\null\hfill{\small\textbf{MESSUNGSFEHLER.}}\hfill\makebox[0pt][r]{\fontsize{13}{16}\selectfont 501}\par\medskip
\begin{center}
\begingroup\renewcommand{\arraystretch}{1.3}\setlength{\tabcolsep}{11pt}
\begin{tabular}{@{\vrule width .8pt\hspace{11pt}}c|c|c@{\hspace{11pt}\vrule width .8pt}}
\noalign{\hrule height .8pt}
[im Jahre & Anzahl der Ables. & mittl. Fehler]\\[4pt]
\hline
1821 & \makebox[2.4em][r]{410} & \makebox[3em][r]{$1\overset{\prime\prime}{,}25$}\\[7pt]
1822 & \makebox[2.4em][r]{717} & \makebox[3em][r]{$1{,}57$}\\
1823 & \makebox[2.4em][r]{449} & \makebox[3em][r]{$1{,}61$}\\
1824 & \makebox[2.4em][r]{843} & \makebox[3em][r]{$1{,}60$}\\
1825 & \makebox[2.4em][r]{609} & \makebox[3em][r]{$1{,}55$}\\
\hline
1822---1825 & \makebox[2.4em][r]{2618} & \makebox[3em][r]{$1{,}58$}\\
\end{tabular}
\endgroup
\end{center}
\par\noindent\hspace*{\parindent}Aus der so sehr nahen Übereinstimmung der Jahre 1822---1825 lässt sich schon schliessen, dass der abweichende Werth von 1821 nicht zufällig ist; in der That waren die Striche an \textsc{Schumacher}s Theodolithen noch etwas schöner, als an meinem eigenen. Jener rührte noch aus der Zeit her, wo \textsc{Reichenbach} selbst die Werkstatt hatte; der meinige war bei \textsc{Reichenbach} bestellt, aber doch unter seiner Aufsicht von \textsc{Ertel} gearbeitet. Auch an meinem 8-zölligen (von 1829) ist die Theilung merklich weniger schön, als an dem 8-zölligen von \textsc{Reichenbach}, den ich 1813 acquirirte; jener hat aber ein \textls[120]{viel} besseres .Fernrohr.

\par\noindent\hspace*{\parindent}Ich sehe, dass ich Ihnen nun doch das Wesentlichste über diesen Gegenstand geschrieben habe, und bemerke also nur noch der Vollständigkeit wegen eine wesentliche Abkürzung jener Rechnung. Es ist nemlich $\Sigma l=n\lambda$, also
\[
\Sigma(l-\lambda)^2=\Sigma ll-2n\lambda\lambda+n\lambda\lambda=\Sigma ll-n\lambda\lambda
\left[=\Sigma ll-\frac{(\Sigma l)^2}{n}\right],
\]
\noindent also
\[
k=\frac{\Sigma ll}{n-1}-\frac{n}{n-1}\lambda\lambda
\left[=\frac{\Sigma ll}{n-1}-\frac{(\Sigma l)^2}{n(n-1)}\right].
\]
\par\bigskip\bigskip
\begin{center}{\large [2.]}\end{center}
\par\noindent\hspace*{\parindent}Zur Ausmittelung der Fehler, die aus der Theilung und dem Ablesen entspringen, wurden alle vollständigen Ablesungen am Theodolithen, 449 an der Zahl [aus dem Jahre 1823], auf folgende Art discutirt.

% Printed 502; principal PDF 509; Google PDF 513.
\par\clearpage
\noindent\makebox[0pt][l]{\fontsize{13}{16}\selectfont 502}\hfill{\small\textbf{NACHLASS.}}\hfill\null\par\medskip
\par\noindent\hspace*{\parindent}Summe der Ablesungen am ersten und dritten Index weniger Summe der Ablesungen am zweiten und vierten $=l$.
\[
\begin{aligned}
\Sigma l&=+133\\
\Sigma ll&=\phantom{+}18521,
\end{aligned}
\]
\noindent also der mittlere Werth von $l=\left[\frac{133}{449}=\right]+0\overset{\prime\prime}{,}30=\lambda$;
\[
\text{mittlerer Werth von }(l-\lambda)^2=
\left[\frac{1}{448}\left(18521-\frac{133^2}{449}\right)=\right]41{,}254.
\]
\par\noindent\hspace*{\parindent}Also der mittlere Fehler bei Einem Vernier inclusive des Theilungsfehlers:
\[
\left[\frac{\sqrt{41{,}254}}{2}=\right]\pm3\overset{\prime\prime}{,}211.
\]
\par\noindent\hspace*{\parindent}An \textsc{Schumacher}s Theodolithen gaben 410 Ablesungen vom Jahre 1821
\[
\begin{aligned}
\Sigma l&=-1896\\
\Sigma ll&=\phantom{-}18982.
\end{aligned}
\]
\par\noindent\hspace*{\parindent}Mittlerer Werth von $l=-4\overset{\prime\prime}{,}62=\lambda$.
\par\noindent\hspace*{\parindent}Mittlerer Werth von $(l-\lambda)^2=24{,}97$.
\par\noindent\hspace*{\parindent}Mittlerer Ablesungsfehler Eines Verniers inclusive des Theilungsfehlers $=\pm2\overset{\prime\prime}{,}499$.
\par\bigskip\bigskip
\par\noindent\hspace*{\parindent}Der Theodolith gab 1822 aus 717 Ablesungen
\[
\begin{aligned}
\Sigma l&=+212\\
\Sigma ll&=\phantom{+}28433.
\end{aligned}
\]
\par\noindent\hspace*{\parindent}Mittlerer Werth von $l=+0\overset{\prime\prime}{,}30=\lambda$.
\par\noindent\hspace*{\parindent}Mittlerer Werth von $(l-\lambda)^2=39{,}623$.
\par\noindent\hspace*{\parindent}Mittlerer Ablesungsfehler bei Einem Vernier inclusive des Theilungsfehlers $=\pm3\overset{\prime\prime}{,}147$.
\par\bigskip\bigskip
\par\noindent\hspace*{\parindent}Dies Resultat harmonirt so nahe mit dem von 1823, dass wir beide zusammenfassen können. Wir haben also aus 1166 Ablesungen:

% Printed 503; principal PDF 510; Google PDF 514.
\par\clearpage
\noindent\null\hfill{\small\textbf{MESSUNGSFEHLER.}}\hfill\makebox[0pt][r]{\fontsize{13}{16}\selectfont 503}\par\medskip
\[
\begin{aligned}
\Sigma l&=+345\\
\Sigma ll&=\phantom{+}46954\\
\lambda&=+0\overset{\prime\prime}{,}296.
\end{aligned}
\]
\par\noindent\hspace*{\parindent}Mittlerer Werth von $(l-\lambda)^2=40{,}216$.
\par\noindent\hspace*{\parindent}Mittlerer Fehler bei Einem Vernier $=\pm3\overset{\prime\prime}{,}1708$.
\par\noindent\hspace*{\parindent}Das Vorkommen der einzelnen Werthe von $l$, ohne Rücksicht auf das Zeichen, war:
\begin{center}
\begingroup\renewcommand{\arraystretch}{1.15}\setlength{\tabcolsep}{10pt}
\begin{tabular}{@{\vrule width .8pt\hspace{10pt}}c|c@{\hspace{10pt}\vrule width .8pt\hspace{10pt}}c|c@{\hspace{10pt}\vrule width .8pt\hspace{10pt}}c|c@{\hspace{10pt}\vrule width .8pt}}
\noalign{\hrule height .8pt}
\makebox[0pt][l]{$[$}\makebox[2em]{$l$} & Anzahl & \makebox[2em]{$l$} & Anzahl & \makebox[2em]{$l$} & Anzahl $]$\\[4pt]
\hline
\makebox[.6em][r]{0}\rlap{$''$} & \makebox[1.8em][r]{45} & \makebox[1.2em][r]{6}\rlap{$''$} & \makebox[1.8em][r]{133} & \makebox[1.2em][r]{11}\rlap{$''$} & \makebox[1.2em][r]{38}\\
\makebox[.6em][r]{1} & \makebox[1.8em][r]{106} & \makebox[1.2em][r]{7} & \makebox[1.8em][r]{104} & \makebox[1.2em][r]{12} & \makebox[1.2em][r]{24}\\
\makebox[.6em][r]{2} & \makebox[1.8em][r]{108} & \makebox[1.2em][r]{8} & \makebox[1.8em][r]{82} & \makebox[1.2em][r]{13} & \makebox[1.2em][r]{15}\\
\makebox[.6em][r]{3} & \makebox[1.8em][r]{123} & \makebox[1.2em][r]{9} & \makebox[1.8em][r]{84} & \makebox[1.2em][r]{14} & \makebox[1.2em][r]{9}\\
\makebox[.6em][r]{4} & \makebox[1.8em][r]{113} & \makebox[1.2em][r]{10} & \makebox[1.8em][r]{52} & \makebox[1.2em][r]{15} & \makebox[1.2em][r]{4}\\
\makebox[.6em][r]{5} & \makebox[1.8em][r]{126} &  &  &  & \\
\end{tabular}
\endgroup
\end{center}
\par\noindent\hspace*{\parindent}Im Jahr 1824 gab der Theodolith aus 843 Ablesungen
\[
\begin{aligned}
\Sigma l&=+104\\
\Sigma ll&=\phantom{+}34506\\
\lambda&=+0\overset{\prime\prime}{,}12.
\end{aligned}
\]
\par\noindent\hspace*{\parindent}Mittlerer Werth von $(l-\lambda)^2=40{,}966$.
\par\noindent\hspace*{\parindent}Mittlerer Fehler bei Einem Vernier inclusive des Theilungsfehlers $=\pm3\overset{\prime\prime}{,}200$.
\par\noindent\hspace*{\parindent}Also aus allen 3 Jahren [1822---1824] 2009 Ablesungen
\[
\begin{aligned}
\Sigma l&=+449\\
\Sigma ll&=\phantom{+}81460\\
\lambda&=+0\overset{\prime\prime}{,}2235.
\end{aligned}
\]
\par\noindent\hspace*{\parindent}Mittlerer Werth von $(l-\lambda)^2=40{,}518$.
\par\noindent\hspace*{\parindent}Mittlerer Fehler Eines Verniers $=\pm3\overset{\prime\prime}{,}1827$. Das Vorkommen der einzelnen Werthe von $l$, ohne Rücksicht auf das Zeichen, war:

% Printed 504; principal PDF 511; Google PDF 515.
\par\clearpage
\noindent\makebox[0pt][l]{\fontsize{13}{16}\selectfont 504}\hfill{\small\textbf{NACHLASS.}}\hfill\null\par\medskip
\begin{center}
\begingroup\renewcommand{\arraystretch}{1.15}\setlength{\tabcolsep}{10pt}
\begin{tabular}{@{\vrule width .8pt\hspace{10pt}}c|c@{\hspace{10pt}\vrule width .8pt\hspace{10pt}}c|c@{\hspace{10pt}\vrule width .8pt\hspace{10pt}}c|c@{\hspace{10pt}\vrule width .8pt}}
\noalign{\hrule height .8pt}
\makebox[0pt][l]{$[$}\makebox[2em]{$l$} & Anzahl & \makebox[2em]{$l$} & Anzahl & \makebox[2em]{$l$} & Anzahl $]$\\[4pt]
\hline
\makebox[.6em][r]{0}\rlap{$''$} & \makebox[1.8em][r]{80} & \makebox[1.2em][r]{7}\rlap{$''$} & \makebox[1.8em][r]{191} & \makebox[1.2em][r]{14}\rlap{$''$} & \makebox[1.2em][r]{13}\\
\makebox[.6em][r]{1} & \makebox[1.8em][r]{182} & \makebox[1.2em][r]{8} & \makebox[1.8em][r]{155} & \makebox[1.2em][r]{15} & \makebox[1.2em][r]{6}\\
\makebox[.6em][r]{2} & \makebox[1.8em][r]{185} & \makebox[1.2em][r]{9} & \makebox[1.8em][r]{138} & \makebox[1.2em][r]{16} & \makebox[1.2em][r]{1}\\
\makebox[.6em][r]{3} & \makebox[1.8em][r]{195} & \makebox[1.2em][r]{10} & \makebox[1.8em][r]{101} & \makebox[1.2em][r]{17} & \makebox[1.2em][r]{1}\\
\makebox[.6em][r]{4} & \makebox[1.8em][r]{205} & \makebox[1.2em][r]{11} & \makebox[1.8em][r]{63} & \makebox[1.2em][r]{18} & \makebox[1.2em][r]{---}\\
\makebox[.6em][r]{5} & \makebox[1.8em][r]{217} & \makebox[1.2em][r]{12} & \makebox[1.8em][r]{38} & \makebox[1.2em][r]{19} & \makebox[1.2em][r]{1}\\
\makebox[.6em][r]{6} & \makebox[1.8em][r]{214} & \makebox[1.2em][r]{13} & \makebox[1.8em][r]{23} &  & \\
\end{tabular}
\endgroup
\end{center}
\begin{center}\rule{3.5cm}{.4pt}\end{center}
\par\bigskip
\par\noindent\hspace*{\parindent}Am [\textsc{Borda}schen] Kreise sind [1823] 46 vollständige Ablesungen gemacht. Es fand sich
\[
\begin{aligned}
\Sigma l&=-174\\
\Sigma ll&=\phantom{-}1534\\
\lambda&=-3\overset{\prime\prime}{,}78.
\end{aligned}
\]
\par\noindent\hspace*{\parindent}Mittlerer Werth von $(l-\lambda)^2=19{,}46$.
\par\noindent\hspace*{\parindent}Mittlerer Ablesungsfehler Eines Verniers inclusive des Theilungsfehlers $=\pm2\overset{\prime\prime}{,}206$.
\par\noindent\hspace*{\parindent}Der Kreis gab 1821 mit 92 Ablesungen
\[
\begin{aligned}
\Sigma l&=-650\\
\Sigma ll&=\phantom{-}11578\\
\lambda&=-7\overset{\prime\prime}{,}07.
\end{aligned}
\]
\par\noindent\hspace*{\parindent}Mittlerer Werth von $(l-\lambda)^2=76{,}767$.
\par\noindent\hspace*{\parindent}Mittlerer Ablesungsfehler Eines Verniers inclusive des Theilungsfehlers $=\pm4\overset{\prime\prime}{,}381$.
\par\noindent\hspace*{\parindent}Am Kreise sind 1822 gemacht 137 Ablesungen; diese gaben
\[
\begin{aligned}
\Sigma l&=-543\\
\Sigma ll&=\phantom{-}6155\\
\lambda&=-3\overset{\prime\prime}{,}96.
\end{aligned}
\]
\par\noindent\hspace*{\parindent}Mittlerer Werth von $(l-\lambda)^2=29{,}432$.

% Printed 505; principal PDF 512; Google PDF 516.
\par\clearpage
\noindent\null\hfill{\small\textbf{MESSUNGSFEHLER.}}\hfill\makebox[0pt][r]{\fontsize{13}{16}\selectfont 505}\par\medskip
\par\noindent\hspace*{\parindent}Mittlerer Ablesungsfehler bei Einem Vernier inclusive des Theilungsfehlers $=\pm2\overset{\prime\prime}{,}713$.
\par\noindent\hspace*{\parindent}Erlaubt man sich, die Messungen von 1822 und 1823 zu verbinden, so wird [bei 183 Ablesungen]
\[
\begin{aligned}
\Sigma l&=-717\\
\Sigma ll&=\phantom{-}7689\\
\lambda&=-3\overset{\prime\prime}{,}92.
\end{aligned}
\]
\par\noindent\hspace*{\parindent}Mittlerer Werth von $(l-\lambda)^2=\left[26{,}812.\right.$
\par\noindent\hspace*{\parindent}Mittlerer Ablesungsfehler bei Einem Vernier inclusive des Theilungsfehlers $=\pm2\overset{\prime\prime}{,}589.]$
\par\bigskip\bigskip
\begin{center}\rule{2.6cm}{.4pt}\end{center}
\vspace{1.5cm}
{\small
\begin{center}BEMERKUNGEN.\end{center}
\par\noindent\hspace*{\parindent}Der vorstehende Brief an \textsc{Gerling}, Art. [1], ist nach dem Original abgedruckt worden; Art. [2] wurde einem Beobachtungsheft für die hannoversche Gradmessung entnommen.
\par\noindent\hfill\textsc{Krüger.}\hspace*{1em}\par}
\vspace{2cm}
\begin{center}\rule{3.5cm}{.4pt}\end{center}

\vfill
\noindent IX.\hfill 64\hspace*{1cm}\par

% Contextual printed 506; principal PDF 513; Google PDF 517: unnumbered blank leaf.
\par\clearpage
\null

% Contextual printed 507; principal PDF 514; Google PDF 518: unnumbered section title.
\par\clearpage
\vspace*{4cm}
\begin{center}
{\fontsize{22}{30}\selectfont
REDUCTION SCHIEFER WINKEL\par
AUF DEN HORIZONT.\par
REPETITIONSBEOBACHTUNGEN.\par}
\end{center}
\vspace{8.2cm}
\noindent\makebox[.86\linewidth][r]{64\textsuperscript{*}}\par
